% ══════════════════════════════════════════════════════════════════════════
% §5  Sparse Systems & The "Triplet" Rule
% ══════════════════════════════════════════════════════════════════════════
\section{Sparse Systems \& The ``Triplet'' Rule}
\label{sec:sparse}

\subsection{Why Sparse?}

Most matrices that arise from discretising PDEs are \emph{sparse}: the
fraction of non-zero entries is $O(1/n)$ or smaller.  Our heat-rod
tridiagonal has $\sim 3n$ non-zeros out of $n^2$ entries.  A 2-D finite
element mesh on $n$ nodes typically produces $\sim 7n$ non-zeros
(depending on element connectivity).

\begin{center}
\renewcommand{\arraystretch}{1.3}
\begin{tabular}{@{}lcc@{}}
\toprule
& \textbf{Dense} & \textbf{Sparse (CSR)} \\
\midrule
Storage for $10^6 \times 10^6$ & $8\;\text{TB}$ & $\sim 24\;\text{MB}$ \\
Matrix--vector multiply & $O(n^2)$ & $O(\text{nnz})$ \\
Direct solve & $O(n^3)$ & depends on fill-in \\
\bottomrule
\end{tabular}
\end{center}

\subsection{The CSR Format}

\textbf{Compressed Sparse Row (CSR)} stores three arrays:
\begin{itemize}[nosep]
    \item \texttt{values[]}: all non-zero entries, row by row.
    \item \texttt{col\_indices[]}: the column index of each entry.
    \item \texttt{row\_ptr[]}: for each row~$i$, the index into
          \texttt{values} where row~$i$ starts.
\end{itemize}

\begin{example}
For the $3\times 3$ tridiagonal matrix from \S1:
\[
    A = \begin{bmatrix} 2 & -1 & 0 \\ -1 & 2 & -1 \\ 0 & -1 & 2 \end{bmatrix}
\]
\begin{center}
\texttt{values = [2, -1, -1, 2, -1, -1, 2]} \\
\texttt{col\_indices = [0, 1, 0, 1, 2, 1, 2]} \\
\texttt{row\_ptr = [0, 2, 5, 7]}
\end{center}
Total storage: 7 doubles + 7 ints + 4 ints = 100 bytes, versus
72 bytes for the full dense matrix.  The savings become dramatic
when $n > 1000$.
\end{example}

\subsection{The Golden Rule of Assembly}

\begin{warningbox}
\textbf{Never} insert entries one at a time into a compressed sparse matrix:
\begin{lstlisting}
// BAD: O(n^2) total cost -- shifts memory on every insert
for (int i = 0; i < n; ++i)
    A.insert(i, i) = 2.0;
\end{lstlisting}
Each \texttt{insert} may need to shift all subsequent entries to make room.
Over $n$ insertions, this is $O(n^2)$.
\end{warningbox}

\textbf{The Triplet Method:} collect all $(i, j, v)$ entries into a flat
\texttt{std::vector<Triplet>}, then compress once.

\begin{lstlisting}
std::vector<Eigen::Triplet<double>> triplets;
triplets.reserve(3 * n);  // pre-allocate

for (int i = 0; i < n; ++i) {
    triplets.push_back({i, i, 2.0});
    if (i > 0)     triplets.push_back({i, i-1, -1.0});
    if (i < n - 1) triplets.push_back({i, i+1, -1.0});
}

Eigen::SparseMatrix<double> A(n, n);
A.setFromTriplets(triplets.begin(), triplets.end());
\end{lstlisting}

\texttt{setFromTriplets} sorts and compresses in $O(\text{nnz}\log\text{nnz})$.
If duplicate $(i,j)$ pairs exist, their values are \emph{summed} by default---this is exactly what finite element assembly needs.

\subsection{The Triplet Memory Trap}

Each \texttt{Eigen::Triplet<double>} stores three values (row, column,
value)---typically 24~bytes.  For a 1-D rod with $n = 10^6$ and $3n$
non-zeros, that is $\sim\!72$~MB just for the temporary vector.
For a 3-D finite element mesh, the triplet vector can easily exceed the
memory of the final compressed sparse matrix itself.

\begin{tipbox}
Two rules to keep triplet memory under control:
\begin{enumerate}[nosep]
    \item \textbf{Always} call \texttt{triplets.reserve(expected\_nnz)}
          before the assembly loop.  Without it, \texttt{std::vector}
          will re-allocate and copy the entire buffer multiple times,
          temporarily tripling your memory footprint.
    \item \textbf{Free immediately} after compression.  Either scope the
          vector in a block, or call
          \texttt{triplets.clear(); triplets.shrink\_to\_fit();}
          right after \texttt{setFromTriplets}.
\end{enumerate}
\end{tipbox}

\subsection{Full Example: 2-D Poisson on a Grid}

\lstinputlisting[caption={Sparse assembly with triplets for a 2-D Poisson problem.},
                 label={lst:sparse2d}]
                {code/sparse_poisson_2d.cpp}
